\subsection{Fields}\label{sec:mathematical_underpinnings:fields}

A field is a triple $(F, \bigcdot, \bigemptydot)$, where $F$ is a non-empty set while $\bigcdot$ and $\circ$ are binary operations that satisfy the following properties:

\begin{itemize}
	\item $(F, \bigcdot)$ is an abelian group;
	\item $(F\textsuperscript{*}, \circ)$ is an abelian group with $F\textsuperscript{*} = F \setminus 0$ --- where $0$ is the identity element for the operation $\bigcdot$.
	\item \textbf{distributivity}: the operation $\circ$ distributes over the operation $\bigcdot$. Formally, for every $x, y, z \in F$, $x \circ (y \bigcdot z) = (x \circ y) \bigcdot (x \circ z)$.
\end{itemize}

Hence, the operations $\bigcdot$ and $\circ$ are compatible --- in the sense that they can both be applied to elements of the set $F$. Also, each operation has its own inverse and identity elements. In simpler words, a field $(F, \bigcdot, \circ)$ is an abelian group which, in addition to the operation $\bigcdot$, also supports a second operation $\circ$ (over $F\textsuperscript{*} = F \setminus 0$).

\remarkBlock{Common names}{Often, the operation $\bigcdot$ is called \textit{addition}, while the operation $\circ$ is called \textit{multiplication}. Consequently, the identity element of $\bigcdot$ is denoted as $0$ and the inverse of an element $a$ as $-a$, while the identity element of $\circ$ is denoted as $1$ --- with $0 \neq 1$ --- and the inverse of an element $a$ as $a\textsuperscript{-1}$. Also, given a field $(F, \bigcdot, \circ)$, the group $(F\textsuperscript{*}, \circ)$ --- that is, the group obtained by considering the set $F \setminus 0$ and the multiplication operation --- is often called the \textit{multiplicative group of nonzero elements} of $F$.}

Most cryptographic algorithms rely on fields with finitely many elements, called \textbf{finite fields} or also \textit{Galois Fields}. In fact, finite fields allow for performing polynomial arithmetic reliably (i.e., with properties guaranteeing the presence of identity and inverse elements). The number of elements in a finite field is called the field's \textbf{order}. A finite field may be either prime or not prime, as explained below.

\paragraph*{Finite Prime Fields.} A finite prime field has order $p$ for some prime number $p$, often denoted as $\mathbb{F}_p = \{0, 1, 2, \ldots, p-1\}$ with addition and multiplication modulo $p$.

\exampleBlock{An example of a finite prime field}{$(\mathbb{F}_7 = \{0, 1, 2, 3, 4, 5, 6 \}, +, \times)$ is a finite prime field. Instead, the whole set of integers $\mathbb{Z}$ is not a field as it does not contain many inverse elements for the multiplication operation $\times$, e.g., $\nexists x \in \mathbb{Z} : x \times 5 = 1$; conversely, such an element exists in $\mathbb{F}_7$ and it is $3$ (in fact, $3 \times 5 = 15 \equiv 1 \pmod{7}$).}

\paragraph*{Finite Non-Prime Fields.} A finite non-prime field has order $p^n$ for some prime number $p$ and integer $n > 1$,\footnote{The value $n$ needs to be greater than $1$, otherwise $\mathbb{F}_{p^1}$ is really just the same as $\mathbb{F}_p$.} often denoted as $\mathbb{F}_{p^n}$. Intuitively, every finite non-prime field $\mathbb{F}_{p^n}$ is somewhat an extension of a smaller prime field $\mathbb{F}_p$ --- this is why finite non-prime fields are also called \textbf{extension fields}.\footnote{More generically, given a base field $K$, an extension field $L$ is any (larger) field that contains the base field $K$ as a subfield. In the finite-field context of cryptography, an extension field is always a finite non-prime field.} A finite non-prime field (or, equivalently, an extension field) $\mathbb{F}_{p^n}$ is built starting from a finite prime field $\mathbb{F}_p$ and then defining a polynomial $f(x)$ with the following properties:
\begin{itemize}
	\item $f(x)$ has degree $n$;
	\item $f(x)$ is irreducible over $\mathbb{F}_p$, that is, $f(x)$ cannot be written as a product of two non-constant polynomials --- e.g., denoted with $h(x)$ and $g(x)$ --- with coefficients in $\mathbb{F}_p$ (\Cref{sec:mathematical_underpinnings:factorization}). Choosing an irreducible $f(x)$ is crucial, otherwise $\mathbb{F}_{p^n}$ would not be a field --- similarly to how choosing a prime $p$ is crucial, otherwise $\mathbb{F}_p$ would not be a field.
	\exampleBlock{Example of reducible and non-reducible polynomials in $\mathbb{F}_5$}{$f(x) = x^2 + 1$ is reducible over $\mathbb{F}_5$: if we consider $h(x) = (x + 2)$ and $g(x) = (x - 2)$, then $h(x) g(x) = (x + 2) (x - 2) = x^2 - 4 \equiv x ^ 2 + 1 \pmod{5} = f(x)$ --- remember that $-4 \equiv 1 \pmod{5}$. Differently, $f(x) = (x^2 + 2)$ is irreducible over $\mathbb{F}_5$.}
\end{itemize}

Notably, the elements of an extension field $\mathbb{F}_{p^n}$ are not simple integers but polynomials of degree $< n$.

\exampleBlock{An example of a finite non-prime field}{Consider the finite prime field $\mathbb{F}_5 = \{0, 1, 2, 3, 4 \}$, the extension field $\mathbb{F}_{25}$, and the arbitrarily chosen polynomial $f(x) = (x^2 + 2)$ --- which has degree $2$ and is irreducible over $\mathbb{F}_5$, as said above. To find the 25 elements of $\mathbb{F}_{25}$, we start from the set of all polynomials in $x$ whose coefficients are in $\mathbb{F}_5$ --- the set of such polynomials is denoted as $\mathbb{F}_5[x]$ and is equal to \{$c_0 + c_1 x + c_2 x^2 + \ldots + c_d x^d : c_i \in \mathbb{F}_5$ \} (were $d$ denotes the degree of the polynomial, whatever it may be --- as long as it is a finite number). Then, we are interested in $\mathbb{F}_5[x]/(x^2 + 2)$, that is, the set of polynomials $\mathbb{F}_5[x]$ modulus $(x^2 + 2)$ --- similarly to when, starting from the set of integers $\mathbb{Z}$, we define a modulus $p$ in modular arithmetic (\Cref{sec:mathematical_underpinnings:modular_arithmetic}).\footnote{This is why $x^2 + 2$ must be irreducible, as it is the modulus of $\mathbb{F}_{25}$.} The set $\mathbb{F}_5[x]/(x^2 + 2)$ is actually isomorphic (i.e., algebraically identical) to $\mathbb{F}_{25}$, so these two notations identify the same mathematical structure. Now, note that the polynomial $(x^2 + 2)$ in $\mathbb{F}_5[x]/(x^2 + 2)$ is actually equal to the polynomial $(0)$, as $(x^2 + 2)$ modulus $(x^2 + 2)$ is obviously $(0)$; hence, $x^2 + 2 = 0$ means that $x^2 = -2$ and, in $\mathbb{F}_5$, it holds that $-2 \equiv 3 \pmod{5}$: ultimately, $x^2 = 3$. This means that, given any polynomial in $\mathbb{F}_5[x]$, we can replace $x^2$ with $3$, so no polynomial in $\mathbb{F}_5[x]/(x^2 + 2)$ has a degree equal or greater than $2$; for instance, $(x^2 + x + 1)$ in $\mathbb{F}_5[x]/(x^2 + 2)$ is actually equal to $(x - 1)$ or, equivalently, to $(x + 4)$. Consequently, every polynomial in $\mathbb{F}_5[x]/(x^2 + 2)$ has the form $(a + b x)$ for $a, b \in \mathbb{F}_5$. Since we have 5 values for $a$ and 5 values for $b$, $\mathbb{F}_5[x]/(x^2 + 2)$ contains exactly $5^2 = 25$ elements. However, it is better to avoid using $x$ as variable, as $x$ does not correspond to an independent variable anymore since $x^2 = 3$ in our context. Therefore, it is better to use another symbol, usually $\alpha$, where $\alpha$ is equal to $x$ but subject to the rule that $a^2 = 3$. Finally, we can enumerate the elements in $\mathbb{F}_5[x]/(x^2 + 2)$: when $b = 0$ and $a$ goes from $0$ to $4$, the elements are the polynomials $(0)$, $(1)$, $(2)$, $(3)$, and $(4)$, when $b = 1$ and $a$ goes from $0$ to $4$, the elements are the polynomials $(\alpha)$, $(1+\alpha)$, $(2+\alpha)$, $(3+\alpha)$, and $(4+\alpha)$, $\ldots$, and when $b = 4$ and $a$ goes from $0$ to $4$, the elements are the polynomials $(4\alpha)$, $(1+4\alpha)$, $(2+4\alpha)$, $(3+4\alpha)$, and $(4+4\alpha)$.}

You may note that there exist multiple extension fields for a finite prime field: after all, it depends on what polynomial is chosen. Still, all of the extension fields of a finite prime field are isomorphic, hence eventually equivalent. Nonetheless, different polynomials may yield different computational costs, better or more secure implementation, or simply be compliant with standards (\Cref{sec:standards_and_certifications}).
